\clearpage
\newpage
\section{Introduction}

Alzheimer's disease is the most common cause of dementia worldwide and is driven by the progressive accumulation of amyloid-$\beta$ plaques and neurofibrillary tau tangles, a pathological process that begins years before clinical symptoms become apparent~\cite{hardy1992amyloid}.
Identifying individuals who harbour this pathology early is a central goal of current research, yet confirmatory assessment still relies on cerebrospinal fluid assays or molecular neuroimaging, procedures that are expensive, invasive, or capacity constrained.
At the syndrome level, primary progressive aphasia (PPA) presents additional diagnostic challenges: three canonical variants are recognised (non-fluent/agrammatic, semantic, and logopenic)~\cite{gorno2011ppa,gorno2015ppa}, but behavioural overlap between them makes clinical differentiation difficult, particularly in early or atypical presentations~\cite{elahi2017clinicopath}.
Both the need for scalable amyloid screening and the challenge of PPA subtyping motivate the development of non-invasive computational tools that can complement expert judgment.

Spontaneous speech is an attractive candidate for this purpose because it simultaneously engages cognitive planning, lexical access, syntactic construction, motor execution, and prosodic control, all of which may be differentially affected by neurodegenerative pathology.
Compared to imaging or fluid biomarkers, voice can be acquired frequently, remotely, and at low marginal cost~\cite{rezaii2025voiceprints}.
Structured elicitation tasks such as the Cookie Theft picture description from the Western Aphasia Battery~\cite{kertesz1982wab} are already part of routine neuropsychological batteries, and recordings from these tasks can be processed computationally without additional clinical instrumentation.

Early computational approaches to cognitive assessment from speech relied on handcrafted linguistic descriptors and shallow acoustic statistics, demonstrating feasibility but limited generalisation~\cite{fraser2016linguistic}.
The shift toward deep learning, particularly self-supervised speech encoders such as wav2vec~2.0~\cite{baevski2020wav2vec} and contextual text models from the BERT family~\cite{devlin2019bert}, has substantially improved representational capacity.
Community benchmarks such as the ADReSSo challenge~\cite{luz2021adresso} have enabled systematic comparison under controlled conditions, and surveys of this literature consistently report that multimodal systems combining acoustic and textual information outperform either modality alone~\cite{ding2024speechsurvey}.

Despite this progress, most published pipelines report end-to-end accuracy figures that conflate several design choices simultaneously: the modality mix, the fusion architecture, the strategy for aligning acoustic frames with text tokens, and the configuration of the self-supervised encoder.
Because these choices are typically entangled within a single system, it is difficult to determine which technique drives the observed performance and which is incidental to the particular implementation or dataset.
Furthermore, the majority of benchmarked systems target English corpora, and it is unclear how well observed trends transfer to languages with different morphology and prosodic norms.
Spanish-speaking clinical cohorts with pathology-confirmed labels exist through initiatives like the Sant Pau Initiative on Neurodegeneration (SPIN)~\cite{alcolea2019spin}, but these have received limited attention in the multimodal deep learning literature.

\subsection{Motivation and objectives}

This thesis is motivated by the observation that, although multimodal speech-based systems for cognitive assessment have shown promising results, the field lacks a systematic understanding of which individual design choices drive performance.
Published systems vary simultaneously in modality, fusion mechanism, audio-text alignment strategy, and encoder configuration, making it impossible to attribute gains to any single technique.
A structured ablation programme, in which each axis is varied while the others are held at controlled defaults, is needed to disentangle these contributions and guide future system design.

Preliminary investigations into deep learning approaches for AD and PPA classification on the ADReSSo and SPIN cohorts were reported in a companion study~\cite{esteve2025speechadppa}; the present thesis extends that work by decomposing the pipeline into individually ablatable design choices.
The main objectives of this thesis are:

\begin{enumerate}
    \item Quantify the individual and combined contribution of audio, text, and handcrafted acoustic features across three clinical speech tasks: AD classification (ADReSSo, English), amyloid-$\beta$ status prediction (SPIN, Spanish), and PPA variant classification (WAB, Spanish).
    \item Systematically compare seven sequence-to-sequence fusion architectures (self-attention, multi-head self-attention, transformer with positional encoding, cross-attention, gated cross-attention, bidirectional cross-attention, and Mamba), evaluated on all three tasks and replicated across two independent implementations.
    \item Propose and evaluate two novel token-to-token audio-text alignment schemes (uniform and proportional) that operate at the subword level, benchmarked against the word-to-token alignment of CogniAlign~\cite{ortiz2025cognialign}.
    \item Evaluate final-layer versus weighted multi-layer Wav2Vec pooling strategies, replicated across both codebases.
    \item Synthesise the experimental evidence into a recommended hybrid architecture that combines the strongest techniques identified along each ablation axis.
\end{enumerate}

\subsection{Work plan}
\label{ssec:workplan}

% TODO: Update gantt_diagrama.tex with actual project phases and dates.
% The current file contains template data from 2019–2020 that must be replaced
% with the real thesis timeline (literature review, implementation, experiments,
% writing phases, milestones, and defence date).

\begin{figure}[H]
    \centering
    \resizebox{\textwidth}{!}{%
\begin{ganttchart}[y unit title=0.4cm,
y unit chart=0.5cm,
vgrid,hgrid,
title height=1,
today=19,%
today offset=.5,%
today label=Now,%
bar/.style={draw,fill=cyan},
bar incomplete/.append style={fill=yellow!50},
bar height=0.7]{1}{20}

 % months (Feb--Jun 2026), four weeks each
 \gantttitle{Master Thesis (February--June 2026)}{20} \\
 \gantttitle{February}{4}
 \gantttitle{March}{4}
 \gantttitle{April}{4}
 \gantttitle{May}{4}
 \gantttitle{June}{4} \\

 % --- Background & data ---
 \ganttgroup[inline=false]{Background \& data}{1}{8}\\
 \ganttbar[progress=100,name=lit]{Literature review}{1}{5} \\
 \ganttbar[progress=100,name=prep]{Data preprocessing \& embedding cache}{2}{8} \\

 % --- Modelling & experiments ---
 \ganttgroup[inline=false]{Modelling \& experiments}{5}{16}\\
 \ganttbar[progress=100,name=prim]{Primitive stacks (CogniAligned, SER)}{5}{10} \\
 \ganttbar[progress=100,name=abl]{Fusion / pooling / alignment ablations}{8}{12} \\
 \ganttbar[progress=100,name=res]{Unified \texttt{res\_model} implementation}{11}{15} \\
 \ganttbar[progress=90,name=mix]{Layer-mixer \& encoder ablations}{13}{16} \\

 % --- Synthesis & writing ---
 \ganttgroup[inline=false]{Synthesis \& writing}{14}{20}\\
 \ganttbar[progress=85,name=ana]{Results analysis}{15}{18} \\
 \ganttbar[progress=65,name=wri]{Thesis writing}{14}{20} \\
 \ganttbar[progress=0,name=def]{Defence preparation}{19}{20} \\

 % dependencies
 \ganttlink{prep}{prim}
 \ganttlink{prim}{abl}
 \ganttlink{abl}{res}
 \ganttlink{res}{mix}
 \ganttlink{mix}{ana}
 \ganttlink{ana}{wri}
 \ganttlink{wri}{def}

\end{ganttchart}%
}

    \caption{Gantt diagram of the project.}
    \label{fig:gantt}
\end{figure}
